
\documentclass[11pt,a4paper]{article}
\usepackage[utf8]{inputenc}
\usepackage[T1]{fontenc}
\usepackage[french]{babel}
\usepackage{lmodern}
\usepackage{microtype}
\usepackage[a4paper,margin=1.65cm,headheight=15pt]{geometry}
\usepackage{amsmath,amssymb,mathtools}
\usepackage{array,tabularx,multicol,booktabs}
\usepackage{enumitem}
\usepackage[most]{tcolorbox}
\usepackage{xcolor}
\usepackage{tikz}
\usetikzlibrary{arrows.meta,calc,positioning}
\usepackage{pdflscape}
\usepackage{fancyhdr}
\usepackage{hyperref}
\usepackage{needspace}
\usepackage{pagecolor}

\definecolor{fondpage}{RGB}{247,248,250}
\definecolor{encre}{RGB}{31,35,40}
\definecolor{bleu}{RGB}{40,91,135}
\definecolor{vert}{RGB}{55,125,92}
\definecolor{orange}{RGB}{178,105,42}
\definecolor{violet}{RGB}{101,77,145}
\definecolor{rouge}{RGB}{170,72,72}
\definecolor{gris}{RGB}{86,91,99}
\definecolor{bleufonce}{RGB}{28,55,120}
\definecolor{bleuclair}{RGB}{238,243,255}
\definecolor{grisclair}{RGB}{245,245,245}
\definecolor{tableline}{RGB}{45,45,45}
\definecolor{softgray}{RGB}{246,247,249}
\definecolor{deepblue}{RGB}{25,62,115}
\definecolor{accent}{RGB}{20,82,130}

\hypersetup{colorlinks=true,linkcolor=bleufonce,urlcolor=bleufonce,
pdftitle={Parcours de revision -- Dérivation -- Premiere specialite},pdfauthor={}}

\setlength{\parindent}{0pt}
\setlength{\parskip}{0.25em}
\renewcommand{\arraystretch}{1.13}
\setlength{\tabcolsep}{4pt}
\setlist[itemize]{leftmargin=1.25em,itemsep=0.15em,topsep=0.2em}
\setlist[enumerate,1]{label=\textbf{\arabic*.}, leftmargin=1.05cm, itemsep=0.35em, topsep=0.2em}
\setlist[enumerate,2]{label=\textbf{\alph*.}, leftmargin=0.85cm, itemsep=0.25em, topsep=0.15em}

\pagestyle{fancy}
\fancyhf{}
\cfoot{\thepage}

\newcommand{\R}{\mathbb{R}}
\newcommand{\N}{\mathbb{N}}
\newcommand{\sourceq}[1]{ {\scriptsize\textcolor{bleufonce}{(site Première q#1)}}}
\newcommand{\code}[1]{\texttt{#1}}

\newcounter{question}
\newcounter{reponse}

\newenvironment{blocquestions}[1]{%
  \begin{tcolorbox}[enhanced,breakable,colback=bleuclair,colframe=bleufonce,boxrule=0.6pt,arc=2mm,left=2mm,right=2mm,top=2mm,bottom=1.5mm,title={\bfseries #1},coltitle=white,colbacktitle=bleufonce,before skip=3mm,after skip=3mm, fontupper=\large]
  \large
  \begin{enumerate}[leftmargin=*,label=\textbf{\arabic*.},itemsep=0.82em,topsep=0.5em]
  \setcounter{enumi}{\value{question}}
}{%
  \setcounter{question}{\value{enumi}}
  \end{enumerate}
  \end{tcolorbox}
}

\newenvironment{blocreponses}[1]{%
  \begin{tcolorbox}[enhanced,breakable,colback=bleuclair,colframe=bleufonce,boxrule=0.6pt,arc=2mm,left=2mm,right=2mm,top=2mm,bottom=1.5mm,title={\bfseries #1},coltitle=white,colbacktitle=bleufonce,before skip=3mm,after skip=3mm, fontupper=\large]
  \large
  \begin{enumerate}[leftmargin=*,label=\textbf{\arabic*.},itemsep=0.42em,topsep=0.5em]
  \setcounter{enumi}{\value{reponse}}
}{%
  \setcounter{reponse}{\value{enumi}}
  \end{enumerate}
  \end{tcolorbox}
}

\newtcolorbox{correctionbox}[1]{enhanced,breakable,colback=grisclair,colframe=black!55,boxrule=0.55pt,arc=2mm,left=2mm,right=2mm,top=2mm,bottom=1.5mm,title=\textbf{#1},coltitle=black,colbacktitle=black!12,before skip=2mm,after skip=2mm}

\newtcolorbox{methodebox}[1]{enhanced,breakable,colback=white,colframe=bleu!70!black,boxrule=0.55pt,arc=2mm,left=2mm,right=2mm,top=2mm,bottom=1.5mm,title=\textbf{#1},coltitle=white,colbacktitle=bleu!70!black,before skip=2mm,after skip=2mm}

\newcommand{\titreSujet}[2]{%
  \subsection*{#1}
  \addcontentsline{toc}{subsection}{#1}
  \textit{Référence / inspiration : #2}\par\medskip\hrule\medskip
}

\newcommand{\qcm}[4]{%
\begin{center}
\begin{tabularx}{0.96\linewidth}{@{}XX@{}}
\textbf{a.} #1 & \textbf{b.} #2\\[1mm]
\textbf{c.} #3 & \textbf{d.} #4
\end{tabularx}
\end{center}}

\tcbset{enhanced,arc=2mm,outer arc=2mm,boxrule=0.42pt,left=1.15mm,right=1.15mm,top=0.75mm,bottom=0.75mm,boxsep=0.35mm,before skip=0.8mm,after skip=0.8mm,fonttitle=\bfseries\footnotesize,coltitle=encre,before upper=\raggedright\fontsize{7.65}{8.15}\selectfont}
\newtcolorbox{fichebox}[3][]{colback=white,colframe=#2!72!black,colbacktitle=#2!18,coltitle=#2!50!black,borderline west={0.9mm}{0pt}{#2!78!black},title={#3},drop fuzzy shadow=black!5,#1}
\newcommand{\tagcol}[2]{\begin{tcolorbox}[enhanced,colback=#1!11,colframe=#1!45!black,boxrule=0.42pt,arc=2mm,left=1mm,right=1mm,top=0.45mm,bottom=0.45mm,before skip=0mm,after skip=0.85mm,fontupper=\bfseries\scriptsize,colupper=#1!55!black]\centering #2\end{tcolorbox}}
\newtcolorbox{panel}[1]{colback=fondpage,colframe=fondpage,boxrule=0pt,arc=0mm,left=0mm,right=0mm,top=0mm,bottom=0mm,height=168mm,valign=top,before skip=0mm,after skip=0mm}
\newcommand{\titrepage}{\begin{tcolorbox}[colback=white,colframe=bleu!70!black,boxrule=0.65pt,arc=3mm,left=2.5mm,right=2.5mm,top=1.15mm,bottom=1.15mm,before skip=0mm,after skip=1.4mm,drop fuzzy shadow=black!10]
{\Large\bfseries Parcours de révision -- Dérivation}\hfill {\normalsize Première spécialité -- sans calculatrice}
\end{tcolorbox}}

% Tableaux TikZ sobres pour les corrections
\newcommand{\TableauSigneDerivDeux}{%
\begin{center}
\begin{tikzpicture}[x=1cm,y=0.62cm,>=Latex]
\def\L{1.8}\def\W{12.4}\def\H{2.35}
\fill[black!8] (0,1.55) rectangle (\W,\H);
\fill[black!5] (0,0) rectangle (\L,\H);
\draw[tableline, line width=0.45pt] (0,0) rectangle (\W,\H);
\draw[tableline, line width=0.45pt] (\L,0)--(\L,\H);
\draw[tableline, line width=0.45pt] (0,1.55)--(\W,1.55);
\draw[tableline, line width=0.35pt] (5.2,0)--(5.2,\H);
\draw[tableline, line width=0.35pt] (8.6,0)--(8.6,\H);
\node at (0.9,1.95) {$x$};
\node at (0.9,0.75) {$f'(x)$};
\node at (2.25,1.95) {$-\infty$};
\node at (5.2,1.95) {$-3$};
\node at (8.6,1.95) {$1$};
\node at (12.0,1.95) {$+\infty$};
\node at (3.55,0.75) {$+$};
\node at (5.2,0.75) {$0$};
\node at (6.9,0.75) {$-$};
\node at (8.6,0.75) {$0$};
\node at (10.35,0.75) {$+$};
\end{tikzpicture}
\end{center}}

\newcommand{\TableauVariationDeux}{%
\begin{center}
\begin{tikzpicture}[x=1cm,y=0.62cm,>=Latex]
\def\L{1.8}\def\W{12.4}\def\H{3.65}
\fill[black!8] (0,2.82) rectangle (\W,\H);
\fill[black!5] (0,0) rectangle (\L,\H);
\draw[tableline, line width=0.45pt] (0,0) rectangle (\W,\H);
\draw[tableline, line width=0.45pt] (\L,0)--(\L,\H);
\draw[tableline, line width=0.45pt] (0,2.82)--(\W,2.82);
\draw[tableline, line width=0.45pt] (0,1.82)--(\W,1.82);
\draw[tableline, line width=0.35pt] (5.2,0)--(5.2,\H);
\draw[tableline, line width=0.35pt] (8.6,0)--(8.6,\H);
\node at (0.9,3.22) {$x$};
\node at (0.9,2.32) {$f'(x)$};
\node at (0.9,0.9) {$f$};
\node at (2.25,3.22) {$-\infty$};
\node at (5.2,3.22) {$-3$};
\node at (8.6,3.22) {$1$};
\node at (12.0,3.22) {$+\infty$};
\node at (3.55,2.32) {$+$};\node at (5.2,2.32) {$0$};\node at (6.9,2.32) {$-$};\node at (8.6,2.32) {$0$};\node at (10.35,2.32) {$+$};
\node at (2.25,0.35) {};\node at (5.2,1.45) {$29$};\node at (8.6,0.35) {$-3$};\node at (12.0,1.45) {};
\draw[->,line width=0.65pt] (2.35,0.35)--(4.8,1.35);
\draw[->,line width=0.65pt] (5.6,1.35)--(8.2,0.45);
\draw[->,line width=0.65pt] (9.0,0.45)--(11.55,1.35);
\end{tikzpicture}
\end{center}}

\newcommand{\TableauSigneAire}{%
\begin{center}
\begin{tikzpicture}[x=1cm,y=0.62cm,>=Latex]
\def\L{1.8}\def\W{10.8}\def\H{2.35}
\fill[black!8] (0,1.55) rectangle (\W,\H);
\fill[black!5] (0,0) rectangle (\L,\H);
\draw[tableline, line width=0.45pt] (0,0) rectangle (\W,\H);
\draw[tableline, line width=0.45pt] (\L,0)--(\L,\H);
\draw[tableline, line width=0.45pt] (0,1.55)--(\W,1.55);
\draw[tableline, line width=0.35pt] (6.3,0)--(6.3,\H);
\node at (0.9,1.95) {$x$};
\node at (0.9,0.75) {$A'(x)$};
\node at (2.25,1.95) {$0$};
\node at (6.3,1.95) {$3$};
\node at (10.35,1.95) {$6$};
\node at (4.25,0.75) {$+$};
\node at (6.3,0.75) {$0$};
\node at (8.25,0.75) {$-$};
\end{tikzpicture}
\end{center}}

\newcommand{\TableauVariationAire}{%
\begin{center}
\begin{tikzpicture}[x=1cm,y=0.62cm,>=Latex]
\def\L{1.8}\def\W{10.8}\def\H{3.65}
\fill[black!8] (0,2.82) rectangle (\W,\H);
\fill[black!5] (0,0) rectangle (\L,\H);
\draw[tableline, line width=0.45pt] (0,0) rectangle (\W,\H);
\draw[tableline, line width=0.45pt] (\L,0)--(\L,\H);
\draw[tableline, line width=0.45pt] (0,2.82)--(\W,2.82);
\draw[tableline, line width=0.45pt] (0,1.82)--(\W,1.82);
\draw[tableline, line width=0.35pt] (6.3,0)--(6.3,\H);
\node at (0.9,3.22) {$x$};
\node at (0.9,2.32) {$A'(x)$};
\node at (0.9,0.9) {$A$};
\node at (2.25,3.22) {$0$};
\node at (6.3,3.22) {$3$};
\node at (10.35,3.22) {$6$};
\node at (4.25,2.32) {$+$};\node at (6.3,2.32) {$0$};\node at (8.25,2.32) {$-$};
\node at (2.25,0.35) {$0$};\node at (6.3,1.45) {$18$};\node at (10.35,0.35) {$0$};
\draw[->,line width=0.65pt] (2.75,0.45)--(5.8,1.35);
\draw[->,line width=0.65pt] (6.8,1.35)--(9.9,0.45);
\end{tikzpicture}
\end{center}}


\newcommand{\TableauSigneSujetUn}{%
\begin{center}
\begin{tikzpicture}[x=1cm,y=0.62cm,>=Latex]
\def\L{1.8}\def\W{10.8}\def\H{2.35}
\fill[black!8] (0,1.55) rectangle (\W,\H);
\fill[black!5] (0,0) rectangle (\L,\H);
\draw[tableline, line width=0.45pt] (0,0) rectangle (\W,\H);
\draw[tableline, line width=0.45pt] (\L,0)--(\L,\H);
\draw[tableline, line width=0.45pt] (0,1.55)--(\W,1.55);
\draw[tableline, line width=0.35pt] (4.8,0)--(4.8,\H);
\draw[tableline, line width=0.35pt] (7.8,0)--(7.8,\H);
\node at (0.9,1.95) {$x$};
\node at (0.9,0.75) {$f'(x)$};
\node at (2.1,1.95) {$-1$};
\node at (4.8,1.95) {$1$};
\node at (7.8,1.95) {$4$};
\node at (10.45,1.95) {$5$};
\node at (3.35,0.75) {$-$};
\node at (4.8,0.75) {$0$};
\node at (6.3,0.75) {$+$};
\node at (7.8,0.75) {$0$};
\node at (9.2,0.75) {$-$};
\end{tikzpicture}
\end{center}}

\newcommand{\TableauVariationSujetUn}{%
\begin{center}
\begin{tikzpicture}[x=1cm,y=0.62cm,>=Latex]
\def\L{1.8}\def\W{10.8}\def\H{3.65}
\fill[black!8] (0,2.82) rectangle (\W,\H);
\fill[black!5] (0,0) rectangle (\L,\H);
\draw[tableline, line width=0.45pt] (0,0) rectangle (\W,\H);
\draw[tableline, line width=0.45pt] (\L,0)--(\L,\H);
\draw[tableline, line width=0.45pt] (0,2.82)--(\W,2.82);
\draw[tableline, line width=0.45pt] (0,1.82)--(\W,1.82);
\draw[tableline, line width=0.35pt] (4.8,0)--(4.8,\H);
\draw[tableline, line width=0.35pt] (7.8,0)--(7.8,\H);
\node at (0.9,3.22) {$x$};
\node at (0.9,2.32) {$f'(x)$};
\node at (0.9,0.9) {$f$};
\node at (2.1,3.22) {$-1$};
\node at (4.8,3.22) {$1$};
\node at (7.8,3.22) {$4$};
\node at (10.45,3.22) {$5$};
\node at (3.35,2.32) {$-$};\node at (4.8,2.32) {$0$};\node at (6.3,2.32) {$+$};\node at (7.8,2.32) {$0$};\node at (9.2,2.32) {$-$};
\node at (4.8,0.35) {$f(1)$};\node at (7.8,1.45) {$f(4)$};
\draw[->,line width=0.65pt] (2.3,1.35)--(4.4,0.45);
\draw[->,line width=0.65pt] (5.2,0.45)--(7.4,1.35);
\draw[->,line width=0.65pt] (8.2,1.35)--(10.0,0.45);
\end{tikzpicture}
\end{center}}

\begin{document}
\begin{center}
{\LARGE\bfseries Parcours de révision -- Dérivation}\\[1mm]
{\large Première générale -- spécialité mathématiques}\\[1mm]
{\normalsize Sans calculatrice}
\end{center}
\vspace{2mm}
\tableofcontents
\clearpage

\phantomsection
\addcontentsline{toc}{section}{Partie 1 -- Récapitulatif de cours}
\newgeometry{margin=6.5mm}
\pagestyle{empty}
\begin{landscape}
\pagecolor{fondpage}\color{encre}
\thispagestyle{empty}

\fontsize{7.85}{8.45}\selectfont

\titrepage

\noindent
\begin{minipage}[t]{0.242\linewidth}
\tagcol{bleu}{1. NOMBRE DÉRIVÉ}
\begin{panel}{bleu}

\begin{fichebox}{bleu}{Taux de variation}
Entre $a$ et $a+h$, le taux de variation de $f$ est :
\[
\frac{f(a+h)-f(a)}{h}.
\]
Il mesure une variation moyenne de la fonction.
\end{fichebox}

\begin{fichebox}{vert}{Nombre dérivé}
Si la limite existe, le nombre dérivé de $f$ en $a$ est :
\[
 f'(a)=\lim_{h\to0}\frac{f(a+h)-f(a)}{h}.
\]
Il donne la pente instantanée au point d'abscisse $a$.
\end{fichebox}

\begin{fichebox}{orange}{Lecture graphique}
$f'(a)$ est le coefficient directeur de la tangente à la courbe de $f$ au point d'abscisse $a$.

Si la tangente monte : $f'(a)>0$.

Si elle descend : $f'(a)<0$.

Si elle est horizontale : $f'(a)=0$.
\end{fichebox}

\begin{fichebox}{rouge}{Vocabulaire}
$f(a)$ est une ordonnée.

$f'(a)$ est un coefficient directeur.

Ne pas confondre la courbe et sa tangente.
\end{fichebox}

\end{panel}
\end{minipage}\hfill
\begin{minipage}[t]{0.242\linewidth}
\tagcol{vert}{2. TANGENTE}
\begin{panel}{vert}

\begin{fichebox}{vert}{Équation de la tangente}
Si $f$ est dérivable en $a$, l'équation de la tangente à la courbe de $f$ au point d'abscisse $a$ est :
\[
 y=f'(a)(x-a)+f(a).
\]
Ici, $f'(a)$ est le coefficient directeur et $A(a;f(a))$ est le point de contact.
\end{fichebox}

\begin{fichebox}{bleu}{Cas particuliers}
Tangente horizontale :
\[
 f'(a)=0.
\]
Tangente montante : coefficient directeur positif.

Tangente descendante : coefficient directeur négatif.
\end{fichebox}

\end{panel}
\end{minipage}\hfill
\begin{minipage}[t]{0.242\linewidth}
\tagcol{orange}{3. CALCULER UNE DÉRIVÉE}
\begin{panel}{orange}

\begin{fichebox}{orange}{Dérivées usuelles}
\[
\begin{array}{c|c}
f(x)&f'(x)\\ \hline
k&0\\
x&1\\
x^2&2x\\
x^3&3x^2\\
x^n&nx^{n-1}\\
e^x&e^x\\
\frac1x&-\frac1{x^2}\\
\sqrt{x}&\frac1{2\sqrt{x}}
\end{array}
\]
Sur les intervalles où les fonctions sont définies.
\end{fichebox}

\begin{fichebox}{bleu}{Opérations}
\[
(u+v)'=u'+v'
\]
\[
(ku)'=ku'
\]
Si $h(x)=g(ax+b)$, alors :
\[
h'(x)=a\,g'(ax+b).
\]
\[
(uv)'=u'v+uv'
\]
\[
\left(\frac{u}{v}\right)'=\frac{u'v-uv'}{v^2}
\quad (v\neq0).
\]
\end{fichebox}

\begin{fichebox}{rouge}{Erreurs fréquentes}
\[
(uv)'\neq u'v'.
\]
\[
\left(\frac{u}{v}\right)'\neq\frac{u'}{v'}.
\]
Toujours vérifier le domaine avant d'utiliser une formule.
\end{fichebox}

\end{panel}
\end{minipage}\hfill
\begin{minipage}[t]{0.242\linewidth}
\tagcol{violet}{4. VARIATIONS}
\begin{panel}{violet}

\begin{fichebox}{violet}{Lien signe / variations}
Sur un intervalle :
\[
\begin{array}{c|c}
f'(x)>0&f\text{ croissante}\\
f'(x)<0&f\text{ décroissante}\\
f'(x)=0&\text{tangente horizontale possible}
\end{array}
\]
\end{fichebox}

\begin{fichebox}{vert}{Méthode : tableau de variations}
\begin{enumerate}[leftmargin=*,itemsep=0.1em]
\item Calculer $f'(x)$.
\item Factoriser si possible.
\item Étudier le signe de $f'(x)$.
\item En déduire les variations de $f$.
\end{enumerate}
\end{fichebox}

\begin{fichebox}{orange}{Extremum}
Si $f'$ change de signe de $+$ vers $-$, alors $f$ admet un maximum local.

Si $f'$ change de signe de $-$ vers $+$, alors $f$ admet un minimum local.
\end{fichebox}

\end{panel}
\end{minipage}

\end{landscape}
\nopagecolor
\restoregeometry
\pagestyle{fancy}
\clearpage
\section*{Partie 2 -- Questions flash}\addcontentsline{toc}{section}{Partie 2 -- Questions flash}
\begin{blocquestions}{Questions flash -- Dérivation}
  \item Que représente graphiquement le nombre dérivé $f'(a)$ ? \sourceq{1}
  \item Quelle est la formule du taux de variation de $f$ entre $a$ et $a+h$ ? \sourceq{2}
  \item Si $f'(a)>0$, comment est la tangente au point d'abscisse $a$ ? \sourceq{3}
  \item Si $f'(a)<0$, comment est la tangente au point d'abscisse $a$ ? \sourceq{4}
  \item Si $f'(a)=0$, que peut-on dire de la tangente ? \sourceq{5}
  \item Que vaut la dérivée d'une fonction constante ? \sourceq{6}
  \item Que vaut la dérivée de $f(x)=x$ ? \sourceq{7}
  \item Que vaut la dérivée de $f(x)=x^2$ ? \sourceq{8}
  \item Que vaut la dérivée de $f(x)=x^3$ ? \sourceq{9}
  \item Que vaut la dérivée de $f(x)=x^n$ ? \sourceq{10}
  \item Que vaut la dérivée de $f(x)=\dfrac1x$ ? \sourceq{11}
  \item Que vaut la dérivée de $f(x)=\sqrt{x}$ ? \sourceq{12}
  \item Quelle est la formule de la dérivée d'une somme ? \sourceq{13}
  \item Quelle est la formule de la dérivée de $ku$ ? \sourceq{14}
  \item Quelle est la formule de la dérivée d'un produit ? \sourceq{15}
  \item Quelle est la formule de la dérivée d'un quotient ? \sourceq{16}
  \item Si $f'(x)>0$ sur un intervalle, que peut-on dire de $f$ ? \sourceq{17}
  \item Si $f'(x)<0$ sur un intervalle, que peut-on dire de $f$ ? \sourceq{18}
  \item Quelle méthode permet d'étudier les variations d'une fonction dérivable ? \sourceq{19}
  \item Que peut indiquer un changement de signe de $f'$ de $+$ vers $-$ ? \sourceq{20}
  \item Que peut indiquer un changement de signe de $f'$ de $-$ vers $+$ ? \sourceq{21}
  \item Quelle est la formule de l'équation de la tangente à la courbe de $f$ au point d'abscisse $a$ ? \sourceq{22}
  \item Dans cette formule, que représente $f'(a)$ ? \sourceq{23}
  \item Si $f(2)=5$ et $f'(2)=3$, écrire l'équation de la tangente au point d'abscisse $2$. \sourceq{24}
  \item Si $h(x)=g(2x-1)$, quelle est l'expression de $h'(x)$ ? \sourceq{25}
  \item Quelle est la dérivée de $f(x)=3x^2-5x+1$ ? \sourceq{26}
  \item Quelle est la dérivée de $f(x)=-2x^3+4x$ ? \sourceq{27}
  \item Quelle est la dérivée de $f(x)=(x+1)(x-2)$ ? \sourceq{28}
  \item Quelle est la dérivée de $f(x)=x^2(3x+1)$ ? \sourceq{29}
  \item Quelle est l'erreur classique avec $(uv)'$ ? \sourceq{30}
  \item Quelle est l'erreur classique avec $(u/v)'$ ? \sourceq{31}
  \item Quelle est la dérivée de $f(x)=e^x$ ? \sourceq{32}
  \item Quelle est la dérivée de $f(x)=e^{2x-1}$ ? \sourceq{33}
  \item Quelle est la dérivée de $f(x)=3e^x-2x$ ? \sourceq{34}
  \item Si $f'(x)=-(x-2)$, quel est le signe de $f'$ pour $x<2$ ? \sourceq{35}
  \item Si $f'(x)=-(x-2)$, quel est le signe de $f'$ pour $x>2$ ? \sourceq{36}
\end{blocquestions}
\clearpage
\section*{Partie 3 -- QCM d'automatismes}\addcontentsline{toc}{section}{Partie 3 -- QCM d'automatismes}
Pour chaque question, une seule réponse est exacte. Aucune calculatrice n'est nécessaire.
\begin{enumerate}
\item Le nombre $f'(a)$ représente :
\qcm{l'ordonnée du point}{le coefficient directeur de la tangente}{l'abscisse du sommet}{l'aire sous la courbe}
\item La dérivée de $f(x)=4x^2-3x+7$ est :
\qcm{$8x-3$}{$4x-3$}{$8x+7$}{$4x^2-3$}
\item Si $f'(x)>0$ sur un intervalle, alors $f$ est :
\qcm{décroissante}{constante}{croissante}{négative}
\item La dérivée de $f(x)=x^3$ est :
\qcm{$x^2$}{$2x$}{$3x^2$}{$3x$}
\item Si une tangente a pour équation $y=-2x+5$, son coefficient directeur est :
\qcm{$5$}{$-2$}{$3$}{$2$}
\item La formule correcte est :
\qcm{$(uv)'=u'v'$}{$(uv)'=u'v+uv'$}{$(u+v)'=uv$}{$(ku)'=u'$}
\item Si $f'$ passe de positif à négatif, on peut avoir :
\qcm{un maximum local}{un minimum local}{une asymptote}{une racine double}
\item L'équation de la tangente à la courbe de $f$ au point d'abscisse $a$ est :
\qcm{$y=f(a)(x-a)+f'(a)$}{$y=f'(a)(x-a)+f(a)$}{$y=f'(x)(a-x)+f(a)$}{$y=f'(a)x+a$}
\end{enumerate}
\clearpage


\section*{Partie 4 -- Sujets type bac / E3C}
\addcontentsline{toc}{section}{Partie 4 -- Sujets type bac / E3C}

\titreSujet{Sujet 1 -- Fonction rationnelle, signe et variations}{inspiré du sujet zéro officiel : signe d'un trinôme, fonction rationnelle, dérivée et variations}
On considère la fonction $g$ définie sur $\R$ par :
\[
  g(x)=x^2-5x+4.
\]
On considère aussi la fonction $f$ définie sur l'intervalle $[1;8]$ par :
\[
  f(x)=x-5+\frac4x.
\]
\begin{enumerate}
  \item Factoriser $g(x)$, puis étudier le signe de $g(x)$ sur $\R$.
  \item Vérifier que, pour tout $x\in[1;8]$, $f(x)=\dfrac{g(x)}{x}$.
  \item En déduire le signe de $f(x)$ sur $[1;8]$.
  \item Calculer $f'(x)$.
  \item Montrer que, pour tout $x\in[1;8]$, $f'(x)=\dfrac{(x-2)(x+2)}{x^2}$.
  \item Étudier le signe de $f'(x)$ sur $[1;8]$.
  \item Dresser le tableau de variations de $f$ sur $[1;8]$.
  \item Déterminer le minimum de $f$ sur $[1;8]$.
  \item Interpréter graphiquement le résultat obtenu à la question 3.
\end{enumerate}

\titreSujet{Sujet 2 -- Fonction avec paramètre et tangente}{inspiré d'un sujet E3C : tangente, paramètres, produit avec exponentielle et variations}
On considère la fonction $f$ définie sur $\R$ par :
\[
  f(x)=(ax+b)e^{-x},
\]
où $a$ et $b$ sont deux réels.

Sa courbe représentative passe par le point $A(0;4)$. La tangente à la courbe au point $A$ passe par le point $B(1;7)$.

On rappelle que $e^0=1$, que $e^x>0$ pour tout réel $x$, et que la dérivée de $e^{-x}$ est $-e^{-x}$.
\begin{enumerate}
  \item En utilisant $f(0)=4$, déterminer $b$.
  \item Calculer le coefficient directeur de la tangente en $A$.
  \item En déduire la valeur de $f'(0)$.
  \item Montrer que, pour tout réel $x$,
  \[
    f'(x)=(a-ax-b)e^{-x}.
  \]
  \item En déduire la valeur de $a$, puis l'expression de $f(x)$.
  \item Étudier le signe de $f'(x)$ sur $\R$.
  \item Dresser le tableau de variations de $f$ sur $\R$.
  \item Déterminer le maximum de $f$ sur $\R$.
  \item Donner une équation de la tangente à la courbe au point $A$.
\end{enumerate}

\clearpage
\titreSujet{Sujet 3 -- Concentration et seuil réglementaire}{inspiré de sujets E3C : dérivée en contexte, maximum et interprétation}
Après injection d'un produit, la concentration dans le sang est modélisée, en mg/L, par la fonction $C$ définie sur $[0;5]$ par :
\[
  C(t)=4t e^{-t}.
\]
On rappelle que $e^t>0$ pour tout réel $t$ et que la dérivée de $e^{-t}$ est $-e^{-t}$.
\begin{enumerate}
  \item Calculer $C(0)$.
  \item Montrer que, pour tout $t\in[0;5]$,
  \[
    C'(t)=4(1-t)e^{-t}.
  \]
  \item Étudier le signe de $C'(t)$ sur $[0;5]$.
  \item Dresser le tableau de variations de $C$ sur $[0;5]$.
  \item Déterminer l'instant où la concentration est maximale.
  \item Donner la valeur exacte de cette concentration maximale.
  \item On admet que $2<e<3$. Montrer que la concentration maximale est supérieure à $1$ mg/L.
  \item La réglementation impose une concentration inférieure ou égale à $1$ mg/L. On admet que $e^3>12$. Montrer qu'au bout de $3$ h, la réglementation est de nouveau respectée.
  \item Interpréter les résultats précédents dans le contexte.
\end{enumerate}

\end{document}
